\documentclass[11pt,a4paper]{article}
\usepackage[T1]{fontenc}
\usepackage[utf8]{inputenc}
\usepackage{mathpazo,amsmath,amssymb,amsthm,mathtools}
\usepackage[margin=25mm]{geometry}
\usepackage{booktabs,longtable,array,enumitem,microtype,xcolor,xurl,hyperref,fancyhdr}
\hypersetup{colorlinks=true,linkcolor=blue!45!black,urlcolor=blue!45!black,citecolor=blue!45!black,
 pdftitle={Projection before Diophantine encoding: a smaller direct quartic for Con(ZFC)},
 pdfauthor={AI-assisted research prepared for Vladimir Reshetnikov}}
\setlength{\headheight}{15pt}
\pagestyle{fancy}\fancyhf{}
\fancyhead[L]{\small Projection before Diophantine encoding}
\fancyhead[R]{\small ProveIt research contribution}
\fancyfoot[C]{\thepage}
\setlength{\parindent}{0pt}\setlength{\parskip}{5pt}
\setlist{nosep,leftmargin=1.6em}\emergencystretch=2em
\newtheorem{theorem}{Theorem}[section]
\newtheorem{lemma}[theorem]{Lemma}
\newtheorem{proposition}[theorem]{Proposition}
\newtheorem{corollary}[theorem]{Corollary}
\theoremstyle{definition}\newtheorem{definition}[theorem]{Definition}
\theoremstyle{remark}\newtheorem{remark}[theorem]{Remark}
\newcommand{\N}{\mathbb N}\newcommand{\Z}{\mathbb Z}\newcommand{\Q}{\mathbb Q}
\newcommand{\Con}{\operatorname{Con}}\newcommand{\Refut}{\operatorname{Ref}}
\newcommand{\Ccal}{\mathcal C_{222}}\newcommand{\SOS}{\operatorname{SOS}}
\newcommand{\Imp}{\mathsf I}\newcommand{\Rel}{\mathsf R}\newcommand{\All}{\mathsf A}
\newcommand{\Eq}{\mathsf E}\newcommand{\Neg}{\mathsf N}\newcommand{\Ex}{\mathsf X}
\newcommand{\Branch}{\mathcal B}\newcommand{\Res}{\mathcal R}
\numberwithin{equation}{section}
\newcommand{\FinalVars}{7{,}617}
\newcommand{\FinalLast}{7{,}616}
\newcommand{\FinalRows}{4{,}379}
\newcommand{\FinalTerms}{148{,}220}
\newcommand{\FinalCoefficient}{8\,974\,652\,400\,619\,654\,748\,160}
\newcommand{\SavedVars}{20{,}299}
\newcommand{\ReductionPercent}{72.71}

\title{\vspace{-10mm}\textbf{Projection before Diophantine Encoding}\\[2mm]
\large A smaller direct quartic whose insolubility is equivalent to\ \(\Con(\mathrm{ZFC})\)}
\author{AI-assisted research prepared for\\Vladimir Reshetnikov's ProveIt project}
\date{24 October 2026\\\small Repository snapshot \texttt{04a1d79222d5}}
\begin{document}
\maketitle
\begin{abstract}
We construct a completely specified polynomial in \FinalVars\ unrestricted
integer variables, of exact total degree four, whose lack of integer
solutions is provably equivalent in ZFC to syntactic consistency of ZFC.
Its complete collected coefficient list and quadratic residuals are supplied.
It has \FinalTerms\ nonzero monomials and largest absolute coefficient
\(\FinalCoefficient\).

The construction starts directly from the twelve numerical C222 proof
rules. It does not simulate a Turing machine or invoke Jones's universal
polynomial. Relative to the supplied direct quartic in 27,916 integer
variables, the new equation removes \SavedVars\ variables
(\ReductionPercent\ percent). The mechanism is to eliminate computed
formula codes before encoding an unbounded family of proof steps, then
pad proofs by a zero-balance truth rule so that proof length bounds the
remaining rule parameters. After the additional coprimality and remainder-bound refinements, the
local family needs eight witness columns rather than twenty-nine. This
cuts the number of exponentiation certificates from 108 to 40. Signed arithmetic and signed
congruence quotients reduce the cost of the final quartic conversion.

This is a substantial simplification of the direct, degree-four branch.
It is not a minimum-variable record: the repository already contains
nine-positive-variable and ten-integer-variable equations of much higher
degree. We give ordinary mathematical proofs, exact data and algebra
checks, and reproducible builders. The logical calibration imports the
repository's C222-to-ZF and ZF-to-ZFC results. A new end-to-end Lean or
encoded object-language ZFC derivation is not supplied.
\end{abstract}
\clearpage
\begingroup\small\setlength{\parskip}{0pt}\tableofcontents\endgroup
\clearpage

\section{Result, scope, and exact equation}

Throughout, \(\N=\{0,1,2,\ldots\}\); the final equation is over
\(\Z\), without sign restrictions on any quantified coordinate.
The phrase ``simplest equation'' has no unique mathematical ordering:
number of variables, degree, coefficient height, arithmetic circuit size,
and description length can move in different directions. We therefore
identify the comparator and give literal measurements.

The starting point is the fully materialized direct quartic in the
incoming package \emph{Direct and centered Diophantine consistency
equations}, hereafter Round12~\cite{round12}. The repository snapshot also
contains later arithmetic reductions not summarized by the historical
collection README. We inspected both that collection and the active
research, including the new thirteen-integer radical-norm package.
The latter is not the source of the equation in this paper.

\begin{table}[ht]\centering
\begin{tabular}{lrrr}\toprule
Construction & Integer variables & Degree & Monomials\\\midrule
Round12 direct baseline &27,916&4&758,275\\
Signed quotients on that baseline &23,048&4&717,327\\
This paper's completed equation &\FinalVars&4&\FinalTerms\\\bottomrule
\end{tabular}
\caption{A comparison within the direct proof-rule construction. All
variables in this table are unrestricted integers.}
\label{tab:comparison}\end{table}

\begin{definition}[The actual polynomial]
Let \(r_1,\ldots,r_{\FinalRows}\) be the literal quadratic polynomials
in \path{equations/residuals_integer.json}, in its stored order. Define
\begin{equation}\boxed{\quad
 P_{\rm direct}(x_0,\ldots,x_{\FinalLast})
   =\sum_{\ell=1}^{\FinalRows}r_\ell(x_0,\ldots,x_{\FinalLast})^2=0.
 \quad}\label{eq:final}
\end{equation}
Its expanded coefficient list is in
\path{equations/equation_integer.json}. A term
\texttt{[c,[i1,...,ik]]} means \(c x_{i_1}\cdots x_{i_k}\);
indices start at zero, repetitions are powers, and an empty index list
denotes the constant term. Each coefficient is an explicit decimal
integer. The file contains neither an unevaluated oracle nor a variable
exponent, factorial, division, square root, or bounded quantifier.
\end{definition}

The sum notation in \eqref{eq:final} is a fixed finite abbreviation.
The JSON is a convenient exact representation of an ordinary polynomial,
just as a long coefficient table would be. The builders specify every
entry; an independent checker reconstructs the collected sum and compares
it exactly. Definitions of intermediate expressions do not add hidden
quantifiers to this displayed variable count. Every retained multiplication
wire and every natural-number representation witness is counted.

\begin{theorem}[Main ordinary theorem]\label{thm:main}
For the particular polynomial \eqref{eq:final}, with the fixed proof
calculus and the inherited logical calibration described below,
\begin{equation}
\mathrm{ZFC}\ \vdash\
\left(\neg\exists\boldsymbol{x}\in\Z^{\FinalVars}
              P_{\rm direct}(\boldsymbol{x})=0\right)
 \ \longleftrightarrow\ \Con(\mathrm{ZFC}).
\label{eq:main}
\end{equation}
The total degree is exactly four.
\end{theorem}

Here \(\Con(\mathrm{ZFC})\) is the usual arithmetical assertion that
there is no finite syntactic ZFC refutation. No consistency assumption is
used to prove the biconditional. Its proof is completed in
Section~\ref{sec:composition}. It neither proves the left side nor proves
consistency. Global minimality and historical priority of the general
compiler ideas are not asserted.

\section{The direct proof calculus and its calibration}
\label{sec:calculus}

The source is a fixed finite list of twelve single-conclusion numerical
rule schemata. Formula constructors are integer polynomials, beginning
with four times the Cantor pairing function:
\begin{align}
\Imp(a,b)&=2(a+b)(a+b+1)+4b,\nonumber\\
\Rel(a,b)&=\Imp(a,b)+1,\qquad
\All(a,b)=\Imp(a,b)+2,\qquad
\Eq(a,b)=\Imp(a,b)+3,\nonumber\\
\Neg(a)&=\Imp(a,1),\qquad
\Ex(a,b)=\Neg(\All(a,\Neg(b))).
\label{eq:constructors}
\end{align}
The truth code is zero and the refutation target is one. The collision
\(\Imp(0,0)=0\) is already covered by the inherited calculus semantics.
The complete rules are specified in Appendix~\ref{app:rules} and the
frozen source data. Their four shared base parameters are
\(a,b,c,d\in\N\). Unused parameters may be zero.

Each rule has two padded premise labels \(p,q\) and one conclusion label
\(v\). An axiom uses \(p=q=0\). Generalization has one zero premise.
The truth rule has \((p,q,v)=(0,0,0)\), so it has no contribution to any
positive-label balance. Only the bundled axiom has an extra side
condition. It is encoded exactly by
\begin{equation}
 \big((c-a)(c-b)\big)^2-(d+1)^2=0,
\label{eq:bundle-side}
\end{equation}
which, for natural parameters, says \(c\ne a\) and \(c\ne b\).
Indeed a nonzero integer product has a positive absolute value, and one
may choose \(d=|(c-a)(c-b)|-1\). Both signs of the differences are allowed.

The logical facts imported from the repository are
\begin{align}
 \Ccal(1)&\longleftrightarrow\Refut_{\rm ZF},\nonumber\\
 \Refut_{\rm ZF}&\longleftrightarrow\Refut_{\rm ZFC}.
\label{eq:calibration}
\end{align}
The first is the mathematical content of
\texttt{KEndpoint.c222\_iff\_zf\_refutation}; the second is
\texttt{ConstructibleChoice.zf\_refutation\_iff\_zfc\_refutation}.
These declarations and their source files are pinned in the archive.
The second equivalence uses the constructible-universe consistency
transport. The first uses the fixed numerical calculus, its completeness
and the set-theoretic seed calibration. Neither equivalence requires the
execution of K215 in this construction~\cite{calculus,choice}.

These are inherited mathematical theorems, not newly proved claims about
the Python compiler. Their ordinary finite coding, completeness and
constructibility arguments can be performed in ZFC. Referring to an
external Lean proposition by itself would not establish the existence of
an encoded object-language ZFC derivation. We keep that distinction
throughout the paper.

\section{Finite proof balances and one obligation accumulator}
\label{sec:flow}

We recall the direct baseline's proof certificate, since the padding
reduction depends on its exact structure. For each rule instance,
associate the vector
\[
 e_v-e_p-e_q,
\]
where \(e_0\) is interpreted as zero and \(e_j\), \(j>0\), is the
unit vector at positive label \(j\).

\begin{lemma}[Single-conclusion balance certificate]\label{lem:flow}
Target one is derivable if and only if a finite multiset of valid rule
instances has balance \(e_1\).
\end{lemma}
\begin{proof}
Unfold a finite derivation into a proof tree. Each internal premise
occurrence cancels the corresponding conclusion occurrence of its
subproof; only the root target remains. Truth leaves contribute zero.

Conversely, take a finite multiset with balance \(e_1\). Let \(D\) be
the least set of labels derivable using just its finitely supported rule
instances, with the zero label regarded as available. Suppose
\(1\notin D\). Sum the balance over the positive labels outside \(D\).
A rule whose conclusion lies outside \(D\) has at least one premise
outside \(D\); otherwise closure would put its conclusion in \(D\).
It therefore contributes at most zero. A rule with conclusion in \(D\)
also contributes at most zero. But the required target vector contributes
one, a contradiction. This argument permits cycles, repeated premises,
and self-loops. It needs nonnegative rule multiplicities and one
conclusion per rule.
\end{proof}

Suppose a list has length \(N\) and all its labels are at most \(h\).
Define
\begin{align}
 A(X)&=\prod_{i<N}(1+v_iX),\nonumber\\
 B(X)&=(1+X)\prod_{i<N}(1+p_iX)(1+q_iX).
\end{align}
Every coefficient of \(A-B\) has absolute value at most
\(2(h+1)^{2N}\). Since \(h+1\le2^h\) for \(h\in\N\), the base
\begin{equation}
 M=M_0+1,\qquad M_0=2^{2hN+1}
\label{eq:heightbase}
\end{equation}
strictly exceeds this coefficient bound.

\begin{lemma}[Integer evaluation recovers the balance]
If \(A(M)=B(M)\) with \(M\) as above, then \(A=B\), and the
positive-label balances equal \(e_1\).
\end{lemma}
\begin{proof}
If a nonzero integer polynomial with coefficient magnitudes below \(M\)
vanishes at \(M\), let \(j\) be its least nonzero coefficient index.
After dividing by \(M^j\) and reducing modulo \(M\), that coefficient
is a nonzero multiple of \(M\), impossible. Hence \(A=B\).
For each positive integer \(r\), compare multiplicities of the distinct
rational root \(-1/r\). The constant factors coming from zero labels
do not contribute. The extra \(1+X\) contributes the target occurrence.
\end{proof}

Use one accumulator \(S_0,\ldots,S_N\) with
\begin{equation}
\begin{split}
 S_0&=M+1,\qquad S_N=1,\\
 S_{i+1}(1+v_iM)&=S_i(1+p_iM)(1+q_iM).
\end{split}\label{eq:accumulator}
\end{equation}
All linear factors are positive. Starting from \(S_0>0\), every
natural accumulator value satisfying an update is positive, and
telescoping gives \(A(M)=B(M)\). Conversely, expand a finite proof tree
by its open obligations: replace a target obligation by its premises.
Take \(S_i\) to be the product of the current obligation weights
\(1+rM\). This gives the required integral accumulator with no rational
intermediate values. Inserting a truth step leaves \(S_i\) unchanged.

\section{Eliminate branch-local data before encoding}
\label{sec:projection}

Round12's local row quantified three labels, four base parameters,
thirteen constructor wires, six beta coordinates, and three label-bound
slacks: twenty-nine natural witness columns. The large cost is that each
column receives its own binomial certificate in the bounded-quantifier
elimination. Eliminating it only at the final quartic stage misses most
of the possible saving.

\begin{lemma}[Projection through a finite disjunction]
Let a finite collection of branches be specified by polynomial equations,
with additional coordinates \(\boldsymbol u\) that are uniquely computed
by a topologically ordered polynomial graph
\(\boldsymbol u=\boldsymbol f_j(\boldsymbol a)\) in branch \(j\).
Assume these graph values lie in the required domain whenever
\(\boldsymbol a\) does. Then
\[
 \exists\boldsymbol u\ \bigvee_j
  \bigl(\boldsymbol u=\boldsymbol f_j(\boldsymbol a)
              \land R_j(\boldsymbol a,\boldsymbol u)=0\bigr)
 \quad\longleftrightarrow\quad
 \bigvee_j R_j(\boldsymbol a,\boldsymbol f_j(\boldsymbol a))=0.
\]
The branches may share the same scratch-coordinate names and have
different graph assignments.
\end{lemma}
\begin{proof}
A satisfying assignment selects one branch and hence its graph values.
Conversely, choose the graph values of a satisfying projected branch.
Other branches need not vanish. A finite product of sums of squares
implements precisely this disjunction over either \(\N\) or \(\Z\).
\end{proof}

In C222 every constructor wire is a polynomial with nonnegative
coefficients. The exposed labels are computed as follows. Axioms have
two zero premises and a conclusion given by their constructor expression;
modus ponens has \((p,q,v)=(a,\Imp(a,b),b)\); generalization has
\((p,q,v)=(b,0,\All(a,b))\); truth has three zero labels. Thus all
sixteen labels and wires can be substituted inside their selected branch.
The bundled side condition \eqref{eq:bundle-side} is retained, including
its natural parameter \(d\). It is not a deterministic constructor.

Write the resulting branch labels as
\(p_j(\boldsymbol a),q_j(\boldsymbol a),v_j(\boldsymbol a)\), where
\(\boldsymbol a=(a,b,c,d)\). They are nonnegative-coefficient
polynomials. Let \(\Gamma_j\) be zero except in the bundle branch, where
it is the square of \eqref{eq:bundle-side}. If \(S,S'\) are the adjacent
accumulator values, set
\begin{equation}
 D_j=\Gamma_j+
 \left[S'(1+v_jM)-S(1+p_jM)(1+q_jM)\right]^2.
\label{eq:branch}
\end{equation}
Each \(D_j\) is nonnegative, and \(\prod_{j=0}^{11}D_j=0\) chooses a
valid rule and its accumulator update. The update must be inside the
branch after label substitution: making a single global substitution
for labels shared by incompatible branches would be incorrect.

This substitution can greatly increase the degree of a local expression.
It does not increase the final degree beyond four, because the local
expression is retained as a shared arithmetic graph and is flattened
only after the bounded family has been eliminated. It exchanges cheap
deterministic multiplication wires for expensive encoded witness columns.

\section{Padding replaces a separate label bound}
\label{sec:padding}

Define the following fixed nonnegative-coefficient polynomial:
\begin{equation}
 H(N)=\sum_{j=0}^{11}
 \bigl(p_j(N,N,N,N)+q_j(N,N,N,N)+v_j(N,N,N,N)\bigr).
\label{eq:H}
\end{equation}
The invalidity of the bundle side condition on the diagonal does not
matter: \(H\) is a majorant for its label polynomial, not a purported
valid diagonal instance. If \(0\le a,b,c,d\le N\), monotonicity gives
\(p_j,q_j,v_j\le H(N)\) for every branch.

\begin{theorem}[Proof-length padding]\label{thm:padding}
C222 derives one if and only if there is an accumulator certificate of
length \(N\) in which every rule's four base parameters are at most
\(N\), and the base is
\begin{equation}
 M_0=2^{2H(N)N+1},\qquad M=M_0+1.
\label{eq:paddedbase}
\end{equation}
\end{theorem}
\begin{proof}
For completeness take a finite proof tree and its finite sequence of
expansion rules. Choose \(N\) at least the sequence length and at least
every base parameter occurring in those rules. This includes the
bundle's \(d\), chosen from its nonzero product. Pad to exactly \(N\)
steps by truth rules with all four parameters zero. Every label is now
at most \(H(N)\). Choose \(M\) using \eqref{eq:paddedbase}, and only
then construct the accumulator from open obligations as in
Section~\ref{sec:flow}. Padding has no effect on balances or products.

For soundness, monotonicity of \(H\) bounds every label. The
coefficient-height argument applies with \(h=H(N)\). The accumulator
therefore gives a balance certificate, and Lemma~\ref{lem:flow} yields a
derivation. A zero-length list is automatically rejected: its endpoints
would require both \(S_0=M+1\) and \(S_0=1\), while \(M\ge3\).
\end{proof}

The order of choices is essential:
\[
 \text{finite proof parameters}\ \longrightarrow\ N
 \ \longrightarrow\ M\ \longrightarrow\ (S_i)
 \ \longrightarrow\ \text{beta codes and their witness bound }U.
\]
There is no circular requirement that \(M\) dominate a bound which
already contains its own accumulator values. The separate label bound
\(m\) and its three per-row slacks disappear. This is why four rule
parameters may use bound \(N\), while six beta coordinates use a
different bound \(U\).

\section{Exactly ten local witness columns}
\label{sec:ten}

Encode the single finite accumulator with the ordinary beta function
\[
 \beta(x,y,i)=x\bmod(1+(i+1)y).
\]
For fixed index \(i\), use the triples
\((S,q_0,s_0)\) and \((S',q_1,s_1)\) to assert
\begin{align}
 x&=q_0(1+(i+1)y)+S,&
 1+(i+1)y&=S+s_0+1,\nonumber\\
 x&=q_1(1+(i+2)y)+S',&
 1+(i+2)y&=S'+s_1+1.
\label{eq:beta}
\end{align}
All six values are natural. These equations express exact remainders,
including the cases \(y=0\), \(S=0\), or \(S'=0\). Their actual
moduli are always positive. Let \(B_\beta\) denote the sum of their
four residual squares. The one local relation is the polynomial equation
\begin{equation}
 \mathcal P(i,\boldsymbol a,S,q_0,s_0,S',q_1,s_1)
 =B_\beta+\prod_{j=0}^{11}D_j=0.
\label{eq:local}
\end{equation}
It has exactly ten local witnesses: four base parameters and six beta
coordinates. The outer parameters are \(N,x,y,M_0\).

The endpoints assert
\(\beta(x,y,0)=M_0+2\) and \(\beta(x,y,N)=1\), consistently with
\(M=M_0+1\). CRT supplies a beta encoding for every finite accumulator:
choose \(y\) larger than all entries and divisible by all positive
integers up to \(N\), then use the pairwise coprime moduli
\(1+(i+1)y\). This choice of \(y\) belongs to the existence proof;
no factorial-divisibility constraint on \(y\) is added to the polynomial.

\section{Bounded-family elimination with two bounds}
\label{sec:crt}

Write \(\mathcal P=e-f\), where \(e,f\) are nonnegative-coefficient
polynomials represented by two natural arithmetic graphs. For example,
\((e-f)(g-h)=(eg+fh)-(eh+fg)\). These graphs are retained for the
majorant argument even though later arithmetic predicates are compiled
with signed intermediate values.

Let the ten bounds be
\[
 \boldsymbol b=(N,N,N,N,U,U,U,U,U,U).
\]
For fixed outer parameters set
\begin{align}
 L&=e(N,\boldsymbol b)+f(N,\boldsymbol b)+N,\nonumber\\
 T&=L!,\qquad d_i=1+(i+1)T,\qquad
 B=\prod_{i<N}d_i.
\label{eq:crt-product}
\end{align}
The precise fixed-parameter arguments of \(e,f\) are those in the
literal source graph. Their occurrence is not an additional tuple of
unknowns.

The bounded family is equivalent to a finite existential certificate
with encoded columns \(a_0,\ldots,a_9\) and an encoded index \(z\):
\begin{align}
 B&\mid (b_j+1)!\binom{a_j}{b_j+1} &&(0\le j<10),\nonumber\\
 B&\mid 1+(z+1)T,\nonumber\\
 e(z,\boldsymbol a)&\equiv f(z,\boldsymbol a)\pmod B.
\label{eq:bounded-cert}
\end{align}
Here the family means \(\forall i<N\;\exists\boldsymbol w\le
\boldsymbol b\;e(i,\boldsymbol w)=f(i,\boldsymbol w)\).

\begin{lemma}[Soundness of the bounded certificate]
Conditions \eqref{eq:bounded-cert} imply the bounded family.
\end{lemma}
\begin{proof}
Fix \(i<N\), and choose a prime \(p\mid d_i\). Since
\(d_i=1+(i+1)L!\), the prime cannot divide \(L!\), so \(p>L\).
In particular \(p>N\), and \(T\) is invertible modulo \(p\).
The encoded-index condition gives \(z\equiv i\pmod p\).

For a column \(a=a_j\), if \(a<b_j+1\), take \(w_j=a\).
Otherwise the product
\((b_j+1)!\binom a{b_j+1}=a(a-1)\cdots(a-b_j)\)
is divisible by \(p\), so some factor gives
\(a\equiv w_j\pmod p\) with \(0\le w_j\le b_j\).
This argument does not require the prime to exceed every coordinate
bound. The encoded evaluation congruence yields
\(e(i,\boldsymbol w)\equiv f(i,\boldsymbol w)\pmod p\).
Nonnegative-coefficient monotonicity bounds both values by their upper
evaluations, hence by \(L<p\). They are equal as integers.
\end{proof}

\begin{lemma}[Completeness of the bounded certificate]
If the bounded family holds, encoded columns and an index satisfying
\eqref{eq:bounded-cert} exist.
\end{lemma}
\begin{proof}
The finitely many moduli \(d_i\) are pairwise coprime. Indeed a common
prime divisor of \(d_i,d_j\) is coprime to \(T\) and divides
\(|i-j|\), which is at most \(N\le L\) and hence divides \(L!\),
a contradiction. Choose natural CRT representatives with
\(a_j\equiv w_{ij}\pmod{d_i}\) and \(z\equiv i\pmod{d_i}\).
Each falling factorial vanishes modulo every \(d_i\): if its upper
argument is too small the integer binomial coefficient is zero, and
otherwise one factor has the required residue. The evaluation and index
conditions hold modulo each factor and therefore modulo \(B\).
For \(N=0\), the empty product is one and the certificate is vacuous;
the separate endpoints still reject that case.
\end{proof}

To choose \(U\) in a completeness construction, first fix the proof,
\(N,M\), the accumulator and its beta encoding. Choose \(U\) above all
six beta coordinates in the finitely many rows. The ten-column family
then holds with its mixed bounds. Recompute \(T,B\) for those bounds;
one must not simply enlarge the bounds in a preexisting CRT certificate.

Only two different bound factorials are needed. The \(U\)-columns share
\((U+1)!\), and the four \(N\)-columns use
\begin{equation}
 (N+1)!=(N+1)N!.
\label{eq:sharedfact}
\end{equation}
The factor \(N!\) is already required to construct \(B\). Thus moving
to mixed bounds adds no factorial certificate.

\section{Final refinement: eight columns and no bound factorials}
\label{sec:eight}

The ten-column construction is a checked intermediate, yielding the
9,476-integer quartic recorded in the supplementary comparison. The
released equation adds two further exact reductions. They remove eleven
power certificates from that intermediate and do not change the proof
calculus or the final degree.

\subsection{Bound beta remainders by the beta scale}

In completeness we may choose the beta scale \(y\) larger than every
accumulator entry before choosing its code \(x\). Thus it is legitimate
to require both adjacent remainders to satisfy \(0\le S,S'\le y\).
For every \(i\ge0\),
\[
 y<1+(i+1)y\quad\hbox{and}\quad y<1+(i+2)y.
\]
These inequalities also hold at \(y=0\). Consequently the two equations
\begin{equation}
 x=q_0(1+(i+1)y)+S,\qquad
 x=q_1(1+(i+2)y)+S'
 \label{eq:beta8}
\end{equation}
already express the exact two remainders when \(S,S'\le y\).
The two remainder-slack columns and their defining rows can be deleted.
The final local bound vector, in the stored coordinate order, is
\begin{equation}
 \boldsymbol b=(N,N,N,N,y,U,y,U),
 \label{eq:bounds8}
\end{equation}
for \((a,b,c,d,S,q_0,S',q_1)\). The local polynomial is obtained from
\eqref{eq:local} by replacing its four beta residuals by the two in
\eqref{eq:beta8}; the branch factors \eqref{eq:branch} are unchanged.

For soundness, the decoded bounds supply the missing strict remainder
conditions. For completeness, choose \(y\) as stated, choose \(x\) by
CRT, and then choose \(U\) above the finitely many quotient values.
No bound on \(x\) or \(U\) is used to choose \(M\). All witness
choices therefore retain their correct dependency order.

\subsection{Cancel the bound factorials using coprimality}

Enlarge the factorial argument to
\begin{equation}
 L=e(N,\boldsymbol b)+f(N,\boldsymbol b)+N+U+y+1,
 \qquad T=L!.
 \label{eq:L8}
\end{equation}
Every entry \(b_j\) of \eqref{eq:bounds8} satisfies \(b_j+1\le L\).
Hence \(F_j=(b_j+1)!\) divides \(T\). For
\(B=\prod_{i<N}(1+(i+1)T)\), each factor is one modulo \(F_j\), so
\begin{equation}
 B\equiv1\pmod{F_j},\qquad \gcd(B,F_j)=1.
 \label{eq:factorialunit}
\end{equation}
It follows that
\begin{equation}
 B\mid F_j\binom{a_j}{b_j+1}
 \quad\longleftrightarrow\quad
 B\mid\binom{a_j}{b_j+1}.
 \label{eq:cancel}
\end{equation}
Thus the final coordinate-range certificates contain just eight binomial
divisibilities, with no separate factorial value for any of their bounds.
The factorials \(L!\) and \(N!\) remain because they are part of the
CRT product construction; \((U+1)!\) disappears entirely.

\begin{theorem}[The final bounded-family interface]\label{thm:eight}
For the eight-column row, fixed outer parameters and bounds
\eqref{eq:bounds8}, take \(T\) from \eqref{eq:L8}. The bounded local
family is equivalent to existence of encoded columns and an encoded
index satisfying
\begin{align}
 B&\mid\binom{a_j}{b_j+1} &&(0\le j<8),\nonumber\\
 B&\mid1+(z+1)T,\qquad
 e(z,\boldsymbol a)\equiv f(z,\boldsymbol a)\pmod B.
 \label{eq:cert8}
\end{align}
\end{theorem}
\begin{proof}
The proof of the bounded-certificate lemmas remains valid when their
factorial argument is enlarged: it is still at least the two upper
evaluations and \(N\), and it still ensures pairwise coprime moduli.
Apply those lemmas to the eight-coordinate source and use
\eqref{eq:factorialunit}--\eqref{eq:cancel} on each column. At \(N=0\),
\(B=1\), so the same argument applies to the empty product; the endpoint
contradiction still excludes that case from the complete equation.
\end{proof}

Two deleted binomial certificates save six powers. Removing the one
bound-factorial certificate saves five more. The final complete ledger
therefore has forty power certificates, eleven binomial certificates,
twenty-seven division/remainder certificates, and two factorial
certificates. The slightly larger factorial argument itself is a fixed
polynomial expression and does not add an arithmetic representation
certificate.

\section{Removing all nonpolynomial notation}
\label{sec:primitives}

Equations \eqref{eq:paddedbase}, \eqref{eq:L8} and \eqref{eq:cert8} are an
explanation of the construction, not the final Diophantine language.
The delivered compiler replaces every power, factorial, binomial and
remainder by fixed finite polynomial certificates. It retains the
baseline's pinned arithmetic interfaces~\cite{arith}. The principal
identities are stated here so that the elimination is reviewable.

\subsection{The finite product}
With \(T\ge1\), let
\[
 Q=(1+(N+1)T)^N+2.
\]
Then the CRT product is represented by
\begin{equation}
 B=\left(T^N N!\binom{Q+N}{N}\right)\bmod(TQ-1).
\label{eq:Bformula}
\end{equation}
The modulus is positive, including \(N=0\), because \(Q\ge3\).
For completeness, this identity has a short direct proof. One has
\[
 T^N N!\binom{Q+N}{N}
 =\prod_{j=1}^N(TQ+jT)
 \equiv\prod_{j=1}^N(1+jT)=B\pmod{TQ-1}.
\]
If \(A=(1+(N+1)T)^N\), then
\(B\le(1+NT)^N\le A<TQ-1\), because \(Q=A+2\) and \(T\ge1\).
Thus the indicated remainder is exactly \(B\). The empty product
case \(N=0\) gives one on both sides.

\subsection{Binomial coefficients and factorials}
For \(n,k\in\N\), put \(q=2^n+1\). The base-\(q\) digit identity
gives
\begin{equation}
 \binom nk=
 \left\lfloor\frac{(q+1)^n}{q^k}\right\rfloor\bmod q.
\label{eq:binom}
\end{equation}
Every binomial coefficient in the expansion of \((q+1)^n\) is below
\(q\), so there is no carry. If \(k>n\), the extracted digit is zero.
For \(r=(n+2)^{n+2}\) and \(C=\binom{r+n}{n}\), the pinned factorial
identity is
\begin{equation}
 n!=\left\lfloor\frac{(r+1)^n+C-1}{C}\right\rfloor.
\label{eq:fact}
\end{equation}
It includes \(n=0\); \(C\) is always positive. To see the identity for
\(n\ge2\), write
\[
 \frac{(r+1)^n}{C}
  =n!\prod_{j=1}^n\frac{r+1}{r+j}.
\]
The product is at most one and is greater than
\(1-n(n-1)/(2r)\). Since \(n!n(n-1)/2<r\), the ratio lies in
\((n!-1,n!]\), so its ceiling is \(n!\). Integer ceiling division
is exactly the right side of \eqref{eq:fact}. The cases zero and one
are immediate. The predecessor \(C-1\) and all integer quotients are
represented by polynomial rows.

\subsection{Pell and exponentiation certificates}
Write
\[
 \mathcal N(a,x,y)=x^2+y^2-a^2y^2-1.
\]
For \(a>1\), the graph of the nonnegative Pell coordinates at index
\(n\) is represented by \(n\le y\) and the disjunction of
\((x,y)=(1,0)\) with existence of natural \(u,w,s,t,b\) satisfying
\begin{gather}
 \mathcal N(a,x,y)=\mathcal N(a,u,w)=\mathcal N(b,s,t)=0,\nonumber\\
 b>1,\quad w>0,\quad y^2\mid w,\nonumber\\
 b\equiv1\pmod{4y},\quad b\equiv a\pmod u,
 \quad s\equiv x\pmod u,\quad t\equiv n\pmod{4y}.
\label{eq:pell}
\end{gather}
These are the finite Matiyasevich conditions used by the frozen baseline.
The compiler includes their zero-index branch. It does not silently
assume that all moduli in an inactive branch are nonzero.

The natural power graph has the explicit cases
\[
 (k=0\land m=1)\ \lor\
 (k>0\land[(n=0\land m=0)\lor(n>0\land\mathcal P(m,n,k))]).
\]
The nonzero case \(\mathcal P\) uses natural \(w,a,t,z,x,y\), the
Pell graph at \((a,k,x,y)\), and
\begin{gather}
 x\equiv y(a\mathbin{\dot-}n)+m\pmod t,
 \qquad 2an=t+n^2+1,\nonumber\\
 m<t,\quad n\le w,\quad k\le w,
 \qquad\mathcal N(w+1,a,wz)=0.
\label{eq:power}
\end{gather}
Truncated subtraction is given its two-branch polynomial definition.
Natural slack variables express inequalities. Sums of squares implement
conjunction and products of nonnegative branch polynomials implement
disjunction. The frozen arithmetic proofs justify these unbounded
interfaces; bounded numerical examples are not used in their place.

\begin{table}[ht]\centering
\begin{tabular}{lrr}\toprule
Arithmetic certificate & Direct baseline & New projected source\\\midrule
Local witness columns &29&8\\
Natural power &108&40\\
Indexed Pell &108&40\\
Binomial &33&11\\
Factorial &3&2\\
Division/remainder &72&27\\\bottomrule
\end{tabular}
\caption{Counts in the complete arithmetic-call ledger, before final
conversion to one quartic over integers.}
\label{tab:calls}\end{table}

\section{Signed arithmetic and congruence quotients}
\label{sec:signed}

The local majorant still uses separate nonnegative expressions \(e,f\),
but there is no reason to require every intermediate arithmetic wire in
the final integer polynomial to be natural. This distinction permits a
more economical final compiler.

\begin{lemma}[A single signed congruence quotient]\label{lem:quotient}
For natural \(a,b,m\), including \(m=0\),
\[
 \exists u,v\in\N\ (a+mu=b+mv)
 \quad\longleftrightarrow\quad
 \exists q\in\Z\ (a-b+mq=0).
\]
\end{lemma}
\begin{proof}
Use \(q=u-v\). Conversely choose
\(u=\max(q,0)\), \(v=\max(-q,0)\). If \(m=0\), both sides say
\(a=b\). Thus the replacement remains valid inside inactive branches
and does not depend on a missing nonzero-modulus hypothesis.
\end{proof}

\begin{lemma}[Division quotients need no independent sign guard]
Suppose \(n,r,s\in\N\), \(q\in\Z\), and
\[
 n=qd+r,\qquad d=r+s+1.
\]
Then \(d>0\), \(0\le r<d\), and \(q\ge0\).
\end{lemma}
\begin{proof}
The first two inequalities follow from the second equation. If
\(q\le-1\), then \(n\le-d+r=-s-1<0\), a contradiction.
\end{proof}

This inference applies to unconditional division graphs whose arguments
have already been proved natural. It does not authorize discarding
arbitrary inequality slacks or positivity guards in conditional
branches. The exact released set of unguarded quotient inputs is recorded
and checked against the call ledger.

With signed arithmetic, the square of a residual is computed directly:
one multiplication for \((a-b)^2\), after an affine subtraction. A pair
representation would compute \(a^2+b^2-2ab\) using three products.
Similarly, a product of two signed values needs one product, whereas
multiplying their positive and negative parts uses four. These are exact
integer identities, so they remain valid at every assignment, not only
on the zero set. Multiplication-gate values may be negative; their sign
is immaterial when the gate equations are squared at the final stage.

\section{One ordinary quartic over unrestricted integers}
\label{sec:quartic}

\begin{lemma}[Signed graph flattening]\label{lem:flatten}
Consider a finite acyclic graph over integer constants, semantic inputs,
addition, subtraction and multiplication, with specified output
equalities. Introduce one integer wire for each genuinely nonlinear
multiplication and substitute affine wires, optionally retaining selected
affine wires to limit expression growth. Then every
retained graph equation is quadratic. Their simultaneous zero set is
equivalent to the original output equalities, with unique reconstructed
graph values for fixed semantic inputs.
\end{lemma}
\begin{proof}
Proceed in topological order. Constants and semantic inputs are affine.
An addition or subtraction of affine forms is affine. A product of two
current affine forms is quadratic; if it remains affine, substitute it,
and otherwise introduce a wire with that quadratic defining row. A retained affine wire simply receives its affine defining
row, with the same topological reconstruction argument.
Consequently all later represented nodes are again affine forms in the
semantic inputs and retained wires. Output equalities give affine rows.
An actual graph evaluation supplies every wire. Conversely, each
defining row forces the corresponding gate value, so induction
reconstructs the whole graph and its outputs over \(\Z\).
\end{proof}

Only semantic inputs whose domain is genuinely natural need a natural
number guard. For each such input \(x\), add three unrestricted
integer witnesses and the quadratic equation
\begin{equation}
 x=u^2+v^2+w(w+1).
\label{eq:natural}
\end{equation}
The right side is nonnegative for every integer \(w\). Conversely every
natural \(x\) has such a representation: \(4x+1\) is not of the
excluded form in the three-square theorem; its three-square
representation has exactly one odd square. Dividing
\(4x+1=(2u)^2+(2v)^2+(2w+1)^2\) by four proves the claim.
This standard conversion is also stated explicitly by Sun~\cite{sun}.

Congruence quotients covered by Lemma~\ref{lem:quotient} have unrestricted
domain. Other unguarded quotient variables have their natural domain
derived by the preceding division lemma, with its hypotheses checked.
The remaining semantic inputs are guarded using \eqref{eq:natural}.
There is no guard for every internal multiplication wire: the signed
graph theorem already accounts for them.

The final polynomial is the sum of the squares of all graph rows,
output rows and semantic guards. Over \(\Z\), a sum of integer squares
vanishes if and only if every row vanishes. It has degree at most four.
The expanded delivered coefficient list contains nonzero monomials of
degree four, which certifies exact degree. In particular, the square of
a genuine natural-number guard contributes a fourth-power term; the
leading homogeneous sum of squares cannot cancel identically over the
reals.

The released dimension count is explicit:
\begin{center}
\begin{tabular}{lr}\toprule
Coordinate role & Count\\\midrule
Semantic inputs &1,268\\
Nonlinear multiplication wires &3,198\\
Retained affine wires &31\\
Natural-number representation witnesses &3,120\\\midrule
Total unrestricted integer coordinates &7,617\\\bottomrule
\end{tabular}
\end{center}
Of the semantic inputs, 201 are signed congruence quotients and 27 are
Euclidean quotients with derived nonnegativity. The remaining 1,040
inputs receive three witnesses each, giving \(3\cdot1040=3120\).
The implementation retains 31 affine wires when substitution would
exceed its fixed sixteen-term expression threshold. Their equations
are linear and are counted explicitly. Together with 110 output
constraints and 1,040 natural guards, the graph gives
\(3198+31+110+1040=4379\) squared residuals.

\section{Proof of the consistency equivalence}
\label{sec:composition}

\begin{proof}[Proof of Theorem~\ref{thm:main}]
Suppose the integer polynomial has a zero. Its sum of squares yields
every quadratic row. The natural guards and derived quotient signs give
the precise semantic domains. Signed graph reconstruction recovers the
finite arithmetic certificates. The pinned Pell, power, binomial,
factorial and remainder theorems therefore give the stated arithmetic
values, without any remaining uninterpreted arithmetic function.

Theorem~\ref{thm:eight} recovers, for each \(i<N\), eight
natural local witnesses bounded by \(\boldsymbol b\). The two beta
equations and decoded remainder bounds identify adjacent values of one
accumulator. Since each
branch polynomial is nonnegative, one branch vanishes and supplies a
valid C222 rule with the correct update. Projection restores all omitted
constructor wires by their polynomial values. The four base parameters
are at most \(N\), so \(H(N)\) bounds every label.
The endpoints, the binary height equation and the accumulator yield
the target balance and hence \(\Ccal(1)\). The two inherited logical
calibrations \eqref{eq:calibration} yield a finite ZFC refutation.

Conversely, suppose there is a finite ZFC refutation. The same
calibration gives a finite C222 derivation. Expand its proof tree, pad it
by truth rules as in Theorem~\ref{thm:padding}, and choose \(N\) before
choosing the base. Define the accumulator, choose a beta encoding and
then a large enough \(U\). Bounded-certificate completeness supplies
the encoded index and columns. The unbounded arithmetic representation
theorems supply all their finite natural witnesses. Replace congruence
quotient pairs by their signed differences. Evaluate the signed graph,
and represent every required natural input by \eqref{eq:natural}.
All final residuals vanish, so the polynomial is zero.

We have proved in ordinary mathematics
\(\exists\boldsymbol x\in\Z^{\FinalVars}P_{\rm direct}(\boldsymbol x)=0
\longleftrightarrow\Refut_{\rm ZFC}\). The finite combinatorics, CRT,
Pell arithmetic and natural-number arguments are formalizable in ZFC;
the inherited logical calibration is likewise an ordinary ZFC
argument. Negating the equivalence gives \eqref{eq:main}. No step assumes
that ZFC is consistent. Exact degree follows from the preceding section.
\end{proof}

The proof has two clearly identified layers. The new contribution proves
the projection, padding, bounded-family specialization, factorial
cancellation and signed
compilation transformations. The arithmetic representation library and
the C222 logical calibration are imported results with frozen source
files and explicit theorem names. The finite verification scripts certify
literal bindings and algebraic identities. They do not independently
establish every unbounded imported theorem, or produce a new kernel proof
of the complete consistency statement.

\section{Verification and reproducibility}
\label{sec:verification}

The archive contains the full expanded polynomial, its quadratic rows,
the final signed arithmetic graph, the projected natural row graph,
the fixed Horn-rule data, builders, independent checks and receipts.
Source hashes pin the reviewed input archive to repository commit
\texttt{04a1d79222d55cabe4ba52897863fc33baa8c268}.
The inherited logical files retain their earlier source pin,
\texttt{9a33342a03d19cf5b061b2083c6b5d3f27f752ed}, as recorded in the
separate provenance manifest.
The native coefficient ports and Turing-machine transition table are not
needed for this direct equation.

The checks cover the following mathematically different obligations:
\begin{enumerate}
\item Exact projection of each of the twelve branches, including all
      constructor definitions, premise/conclusion labels and the retained
      bundle side condition.
\item The two occurrences of the local natural source graph: upper
      evaluation at \((N,\boldsymbol b)\), and encoded evaluation at the
      encoded index and eight encoded columns.
\item The literal mixed bounds, the endpoint equations, the diagonal
      label majorant, and the enlarged factorial argument that permits
      cancellation of every bound factorial.
\item Signed arithmetic identities, congruence replacement and the precise
      semantic-domain list. Inferred nonnegativity is attached to actual
      unconditional division interfaces.
\item Reconstruction of the affine forms and quadratic graph equations,
      followed by exact expansion and collection of their sum of squares.
\item Variable count and liveness, actual degree, monomial count,
      coefficient height, and deterministic regeneration of the final
      artifacts. Mutation controls ensure that selected corrupted inputs
      or artifacts are rejected.
\end{enumerate}

The accompanying verification notes state the exact commands run and
their results. The finite checks establish literal structural identities;
they are not used as a substitute for proofs of the unbounded relations.
No enumeration of all integer tuples, consistency test, proof search for
a ZFC contradiction, or simulation of K215 was performed.

\section{What improved, and what did not}

The principal improvement is a reduction of \(\SavedVars\) unrestricted
integer variables in a completely materialized direct quartic. The
final degree remains four. Coefficient height is a tradeoff: the new
maximum absolute coefficient is \(\FinalCoefficient\), versus
\(320307372009918\) in the baseline; it is not a coefficient-height
improvement. The improvement arises before and during
compilation, not from treating a long list of hidden variables as
notation. The coefficient table is supplied in full.

The most reusable mathematical idea is this ordering of operations:
first substitute deterministic data separately in each finite rule
branch; next use a neutral rule to couple proof length to a parameter
bound; only then encode the bounded family; finally flatten signed
arithmetic. Padding removes a potential circular bound while avoiding
three extra label-slack columns. Bounding remainders by the beta scale
removes two additional columns; factorial cancellation removes one
entire arithmetic certificate.

This equation does not improve the repository's fewest-variable
constructions. Its nine-positive-variable equations and ten-integer
equations occupy a different part of the tradeoff. Sun's ten-unknown
theorem is a current general result~\cite{sun}; no claim that ten is
globally minimal follows from it. Nor does our degree-four construction
prove that a still simpler consistency equation cannot exist.

We investigated a proposed smaller-degree modification of the
thirteen-integer radical construction. Replacing its scale \(Y=qv\)
by a constant multiple of \(v\) appeared to lower polynomial degree,
but destroyed the universal coprimality \(\gcd(H,q)=1\) used by its CRT
recovery. That candidate is not included as a valid equation. It
illustrates why degree counts alone cannot establish a Diophantine
reduction.

Direct relation-algebra encodings are another worthwhile direction.
Formisano, Gavazzi and Omodeo discuss a direct encoding of NBG provability
using an equational relation calculus~\cite{fgo}. The present work uses
the already calibrated numerical C222 source instead of claiming that
their separate axiom-translation project is complete.

\section{Further research questions}

\begin{enumerate}
\item \textbf{Optimize column count before all arithmetic encoding.}
      Can one eliminate one of the four beta remainder/quotient
      columns without introducing an equally expensive bound or division
      certificate elsewhere? A candidate must preserve the shared
      accumulator and both endpoint semantics.
\item \textbf{Share exponentiation certificates.}
      The padded source still uses 40 power certificates. Which repeated
      base/exponent pairs can be shared extensionally, with all zero-base
      and zero-exponent branches preserved? An exact call-graph proof is
      needed; syntactic similarity alone is insufficient.
\item \textbf{Represent the bounded column test more cheaply.}
      Each falling-factorial divisibility condition is currently compiled
      through a binomial coefficient. Is there a lower-cost fixed
      polynomial interface when several columns share the same bound?
\item \textbf{Improve the proof calculus itself.}
      A different equiconsistent finite single-conclusion calculus with
      fewer effective rule parameters could reduce the expensive local
      dimension further. It must include explicit numerical rules and
      two-way syntactic calibration to ZFC refutability.
\item \textbf{Search the coefficient/variable frontier.}
      Affine elimination can increase coefficient height and monomial
      count. Retaining selected affine wires may yield a useful family of
      smaller-coefficient quartics; those tradeoffs should be measured on
      literal polynomials rather than estimated from source syntax.
\item \textbf{Transfer the new source into low-variable compression.}
      The direct quartic can in principle feed a universal compression.
      A worthwhile result must materialize its actual coefficients and
      compare the resulting degree, numeral size and arithmetic graph
      against the existing native constructions. An abstract appeal to
      MRDP is not a completed improvement.
\item \textbf{Formalize the new transformations.}
      The branch projection and padding lemmas are compact targets for
      proof-assistant verification. A complete formal endpoint must also
      identify the emitted coefficient list with its polynomial and
      explicitly internalize the consistency equivalence; an external
      equivalence alone is a different statement.
\end{enumerate}

\appendix
\section{Complete polynomial rule schemata}
\label{app:rules}

Use the constructors in \eqref{eq:constructors}. The eight logical axiom
labels are
\begin{align*}
 \rho_1(a,b)&=\Ex(a,\Eq(a,b)),\\
 \rho_2(a,b,c)&=\Imp(\All(a,\Imp(b,c)),
                     \Imp(\All(a,b),\All(a,c))),\\
 \rho_3(a,b,c)&=\Imp(\Eq(a,b),\Imp(\Rel(c,a),\Rel(c,b))),\\
 \rho_4(a,b,c)&=\Imp(\All(a,\All(b,c)),\All(b,\All(a,c))),\\
 \rho_5(a,b,c)&=\Imp(\Eq(c,a),\Imp(\Rel(c,b),\Rel(a,b))),\\
 \rho_6(a,b)&=\Imp(\Ex(a,\All(a,b)),b),\\
 \rho_G(a,b,c)&=\Imp(\Eq(c,a),
                \Imp(\Eq(c,b),\All(\Rel(a,b),\Eq(a,b)))),\\
 \rho_8(a,b,c,d)&=\Imp(\Imp(\Imp(c,d),a),
                           \Imp(\Imp(a,c),\Imp(b,c))).
\end{align*}
For the bundle define, as nested polynomial abbreviations,
\begin{align*}
 S(a,b,k)&=\All(k,\Imp(\Rel(k,a),\Rel(k,b))),\\
 \operatorname{Req}(b,k,m)
  &=\All(b,\Imp(m,\Neg(\All(b,\Imp(m,\Neg(\Rel(b,k))))))),\\
 \operatorname{Clos}(a,b,k,m)
  &=\All\bigl(b,\Imp(\Rel(b,k),
       \Neg(\All(a,\Imp(m,
          \Neg(\All(a,\Imp(S(a,b,k),\operatorname{Req}(b,k,m))))))))\bigr),\\
 \operatorname{NKS}(a,b,c)
  &=\All(c,\Imp(\operatorname{Clos}(a,b,c,c),
                         \Neg(\operatorname{Req}(b,c,c)))),\\
 \operatorname{VR}(a,b,c)&=\Imp(\Rel(a,b),\All(c,\Rel(a,b))),\\
 \operatorname{Bundle}(a,b,c)
  &=\Neg(\Imp(\operatorname{VR}(a,b,c),\operatorname{NKS}(a,b,c))).
\end{align*}
The literal frozen source and the independent branch checker are the
authoritative finite binding of these expressions to the emitted graph.

\begin{longtable}{llll}
\toprule Rule & Premise 1 & Premise 2 & Conclusion\\\midrule\endhead
Truth &0&0&0\\
Each of \(\rho_1,\ldots,\rho_6,\rho_G,\rho_8\)&0&0&its displayed label\\
Bundle &0&0&\(\operatorname{Bundle}(a,b,c)\)\\
Modus ponens &\(a\)&\(\Imp(a,b)\)&\(b\)\\
Generalization &\(b\)&0&\(\All(a,b)\)\\\bottomrule
\end{longtable}
Only the bundle requires \eqref{eq:bundle-side}; all parameters range over
\(\N\). This is a finite numerical calculus whose formulas happen to
encode logical expressions. The final polynomial does not quantify over
formulas or invoke a parser.

\section{Artifact conventions and integration}

The top-level README specifies the supported build and verification
commands, input provenance, complete file roles and the recommended
equation. The generated metrics file used in this paper is produced from
the same released artifact checked by the verifier. A release hash
manifest covers the final files.

For integration into ProveIt, preserve this contribution as an additional
package under the Diophantine-consistency research material, or retain
the archive in \texttt{docs/incoming}. Its direct compiler does not
replace the low-variable native branch or the inherited formalization.
The frozen source files retain their original licenses and notices;
new files are explicitly identified. No remote repository file was
changed as part of preparing this package.

\begin{thebibliography}{9}
\bibitem{round12} ProveIt research archive,
\emph{Direct and centered Diophantine consistency equations}, Round12,
source package labeled 24 October 2026. Pinned in repository commit
\texttt{04a1d79222d55cabe4ba52897863fc33baa8c268},
\url{https://github.com/VladimirReshetnikov/ProveIt/blob/04a1d79222d55cabe4ba52897863fc33baa8c268/docs/incoming/ConZFC_Direct_Centered_Round12_2026-10-24.zip}.
\bibitem{calculus} ProveIt, numerical C222 calculus and logical endpoint,
\texttt{Code254.lean}, \texttt{KCode222.lean}, \texttt{KEndpoint.lean};
frozen source snapshots and their original repository paths are included
in this package's provenance records. See also
\url{https://github.com/VladimirReshetnikov/ProveIt/tree/9a33342a03d19cf5b061b2083c6b5d3f27f752ed/Computability/BusyBeaver/Lean/BusyBeaver/ZFC}.
\bibitem{choice} ProveIt, \texttt{ConstructibleChoice.lean}, especially
\texttt{zf\_refutation\_iff\_zfc\_refutation}; frozen source snapshot in
the supplied provenance material. The mathematical input is Goedel's
constructible-universe equiconsistency argument.
\bibitem{arith} ProveIt, \texttt{PellRepresentation.lean},
\texttt{PowerRepresentation.lean}, \texttt{FactorialRepresentation.lean},
\texttt{BoundedForallRepresentation.lean}, and
\texttt{BoundedForallCertificate.lean}; pinned source snapshots inherited
from Round12, retained with this contribution.
\bibitem{sun} Z.-W. Sun, \emph{Ten unknowns for Hilbert's tenth problem
over the integers}, arXiv:2609.01594v2, 2026,
\url{https://arxiv.org/html/2609.01594v2}.
\bibitem{fgo} A. Formisano, I. Gavazzi and E. G. Omodeo,
\emph{Diophantine Modeling of Provability in Algebraic Logic},
CILC 2025, CEUR Workshop Proceedings, vol. 4003,
\url{https://ceur-ws.org/Vol-4003/paper13.pdf}.
\end{thebibliography}
\end{document}
